\documentclass[11pt,a4paper]{article}
\usepackage[a4paper,top=15mm,bottom=15mm,left=18mm,right=18mm]{geometry}
\usepackage[T1]{fontenc}
\usepackage{amsmath}
\usepackage{amssymb}
\usepackage{enumitem}

\pagestyle{empty}
\setlength{\parindent}{0pt}
\setlength{\parskip}{4pt}

\begin{document}

\begin{center}
  {\Large\bfseries Answer Key: Reading Scales and Error Intervals}
\end{center}

\section*{Part 1: Reading a scale to the nearest tick}
\begin{enumerate}[label=\textbf{Q\arabic*.}, leftmargin=*, itemsep=2pt]
\item \(14\,\text{cm}\).
\item \(76\,\text{cm}\).
\item One gap \(=20\,\text{ml}\); reading \(=340\,\text{ml}\).
\item One gap \(=2\,\text{ml}\); reading \(=36\,\text{ml}\).
\item One gap \(=10\,\text{g}\); reading \(=440\,\text{g}\).
\item One gap \(=1\,\text{mm}=0.1\,\text{cm}\); reading \(=5.4\,\text{cm}\).
\item One gap \(=0.2\,\text{kg}\); reading \(=23.6\,\text{kg}\).
\item \(1.3\,\text{m}\).
\end{enumerate}

\section*{Part 2: Error intervals}
\begin{enumerate}[label=\textbf{Q\arabic*.}, start=9, leftmargin=*, itemsep=2pt]
\item \(47.5\,\text{cm}\le \text{length}<48.5\,\text{cm}\).
\item \(170\,\text{ml}\le \text{volume}<190\,\text{ml}\).
\item Q1: \(13.5\,\text{cm}\le \text{length}<14.5\,\text{cm}\). Q3: \(330\,\text{ml}\le \text{volume}<350\,\text{ml}\). Q6: \(5.35\,\text{cm}\le \text{length}<5.45\,\text{cm}\). Q7: \(23.5\,\text{kg}\le \text{mass}<23.7\,\text{kg}\).
\item
\begin{itemize}[leftmargin=*, itemsep=2pt]
\item Pencil: half gap \(0.5\,\text{cm}\); \(14.5\,\text{cm}\le L<15.5\,\text{cm}\).
\item Beetle: half gap \(0.5\,\text{mm}\); \(22.5\,\text{mm}\le L<23.5\,\text{mm}\).
\item Flour: half gap \(5\,\text{g}\); \(505\,\text{g}\le m<515\,\text{g}\).
\item Pupil height: half gap \(1\,\text{cm}\); \(147\,\text{cm}\le h<149\,\text{cm}\).
\item Orange juice: half gap \(25\,\text{ml}\); \(225\,\text{ml}\le V<275\,\text{ml}\).
\item Parcel: half gap \(0.1\,\text{kg}\); \(3.3\,\text{kg}\le m<3.5\,\text{kg}\).
\item Medicine: half gap \(0.25\,\text{ml}\); \(2.25\,\text{ml}\le V<2.75\,\text{ml}\).
\item Garden: half gap \(0.05\,\text{m}\); \(12.35\,\text{m}\le L<12.45\,\text{m}\).
\end{itemize}
\end{enumerate}

\section*{Part 3: Using error intervals}
\begin{enumerate}[label=\textbf{Q\arabic*.}, start=13, leftmargin=*, itemsep=2pt]
\item (a) Reading \(31\,\text{cm}\), gap \(1\,\text{cm}\). (b) Reading \(250\,\text{g}\), gap \(10\,\text{g}\). (c) Reading \(1.4\,\text{m}\), gap \(0.1\,\text{m}\) (\(10\,\text{cm}\)).
\item \(22.9\,\text{kg}\le \text{mass}<23.1\,\text{kg}\). No: the interval includes masses above \(23\,\text{kg}\) (e.g. \(23.05\,\text{kg}\)), so Leo cannot be sure.
\item No, Sam is not right. \(6.25\,\text{cm}\le \text{length}<6.35\,\text{cm}\). The true length is not necessarily exactly \(6.3\,\text{cm}\); closer ticks improve precision but do not make it exact.
\item (a) \(195\,\text{g}\le \text{mass}<205\,\text{g}\). (b) \(175\,\text{g}\le \text{mass}<225\,\text{g}\). (c) Use Scales A: its error interval matches the recipe exactly; Scales B could be outside the allowed range.
\item \(295\,\text{ml}\le \text{volume}<305\,\text{ml}\). Total width \(=10\,\text{ml}\), so ticks must be every \(10\,\text{ml}\) (e.g. \(290\,\text{ml},300\,\text{ml},310\,\text{ml}\)).
\item No, not definitely. The shelf and gap both have interval \(79.5\,\text{cm}\le \text{length}<80.5\,\text{cm}\). For example, shelf \(80.4\,\text{cm}\) and gap \(79.6\,\text{cm}\) would not fit. To guarantee a fit, the gap's lower bound would need to exceed the shelf's upper bound.
\item Example: (a) Kitchen measuring jug, ticks every \(50\,\text{ml}\). (b) Reading \(250\,\text{ml}\) gives \(225\,\text{ml}\le \text{volume}<275\,\text{ml}\). (c) Accurate enough for soup or a drink; not accurate enough for medicine.
\end{enumerate}

\end{document}